\documentclass{ltjsbook}
\usepackage{docmute} 
\usepackage{../myStyle} 

\begin{document} 
\VerbatimFootnotes
\chapter{付録；分析環境の準備}
\label{Appendix} 
\section{はじめに}
この授業では統計言語としてのR，それを使う統合開発環境としてのRStudio，そして確率プログラミング言語としてのStanを利用します。大学でのPCルームの利用にはこれらの環境が既に準備されていますが，自分のPCに環境を準備することでいつでも分析ができるようになります。今後の研究室配属やその後の卒業研究などでも活用することになると思いますので，まずは自身のPC環境に準備することを考えてみて下さい。

この付録資料は，これらの環境を準備する方法について解説するものですが，提供されるソフトウェアのバージョンや対応する環境などは日々発展するものですので，必ずしも最適な情報提供になり得ません。基本的な情報は提供いたしますが，執筆時\footnote{この記事は2021年9月14日に執筆されました。}での情報であることも多く，より新しい情報についてはインターネットなどでキーワード検索を行なって，より良いものを探してみてください。

PCの操作が苦手に思っている人のために付言しますが，うまくいかなかったからと言って癇癪を起こしてはいけません。相手は機械ですから，命じられたことを愚直に実行しているだけです。また文房具のようなものですから，恐れて嫌うのではなく，愛して可愛がってあげることが肝要です。もしみなさんの中に，お持ちのPCに名前をつけていない人がいるようでしたら，今すぐ命名することをお勧めします。パソコンが，と思うと腹が立ちますが，わたしのXXXちゃんが，と思うとエラーも「ちょっとご機嫌斜めなのかしら」と受け入れやすくなります\footnote{ちなみに私が初めて手に入れたノートPCには「秀吉」という名前をつけました。Macに変わってからは代々数学者の名前をつけることにしています。私はいま，GaussちゃんやHermiteちゃんを持ち歩いています。}。

また，困った時は教員，TA，先輩などに相談することが重要ですが，その際には症状や経緯を正確に伝えることが必要です。ペットを病院に連れていくときに「何もしてないのに勝手に病気になった。治してほしい」といわれても獣医さんは困ると思います。\textbf{普段の様子や症状についての丁寧な報告}を心がけてください。特にPCは機械ですので，何もしていないのに壊れるということはありません。自分が自覚していないことであっても，「なにかをした」から様子がおかしくなるのです。自分はそんな大それたことはしていない\footnote{例えば計算途中だけど時間が来たので電源をオフにした，と言ったことでも，内部ファイルにアクセしているときの中断であれば，十分にPCを破壊することができます。}，と思っても必ず何かトリガーがあったはずなのです。どこをクリックしたか，どういうコマンドを入れたかという\textbf{履歴をしっかりと把握し，丁寧に報告する}ことを心がけてください。

\section{概略}
環境の導入は，OSの種類によって変わります。WidowsなのかMacなのか，あるいはLinuxなのかといった違いに合わせて導入を考えてください。
OSの種類で言えば，Linuxは全体的にCUI操作が主であり，堅牢かつ合理的な使い方ができるものです。如何せんマイノリティですので，ググってもあまり情報が出てこないところが玉に瑕です。
MacはLinuxベースのOSになっているので，比較的堅牢かつ合理的な使い方ができます。ハードウェアもOSも一つの会社(Apple)が作っているので，サポートもしやすいという利点がありますが，日本では少数派なのが残念なところです。
Windowsは多数派のOSですが，OSメーカはソフトウェア会社でハードウェアが別会社なことが多く，利用される範囲は広いのですが，その分問題が生じたときにケースバイケースになりがちです。多くのハードウェアの違いを吸収するために最大公約数的な設計をするせいか，ソフトウェアデザインの統一性がないところが欠点です。ユーザ数は最も多いので，一生懸命調べたら同じような問題で困っている同じような機体の人に出会える可能性が高いのがせめてもの救いです。また光明としてWindowsOSもLinuxベースの仕組みを取り入れ始め(WLS)，比較的合理的な振る舞いができるようになったことが挙げられますが，OSのグレードによって導入のしやすさが異なります。

ちなみにChromeOSやiOS, Androidなどタブレット・スマートフォンでの統計環境の利用は限界があります。これらは計算結果を利用するフロントエンドとしては非常に高機能なのですが，インタラクティブに使うことには不向きです。いわば書籍のように多くの情報を一方的に与えてはくれるのですが，ユーザが計算するといった働きかけの補助になるような(ノートとペンのような)使い方ができませんので注意してください。

さて，環境の導入には三種類のルートがあります。第一のものが最も基本的なアプローチで，自分の手元の環境を作り上げるというものです。これを\textbf{ローカル環境}での構築と呼びます。ローカル環境は自分だけの環境ということなので，自分の責任でもってしっかりと環境構築をすることができます。第二のアプローチとして，ある程度出来上がった環境を丸ごと取り込む，\textbf{仮想環境}での構築という方法があります。PCのOSの中に別のマシン・別のOSをさらに憑依させて使うようなやりかたで，これができると取り込む環境は完成しているので細かい設定をする必要がありません。問題は「憑依させる」準備がいろいろ大変であるということです。また憑依させたPCはブラウザ経由でアクセスすることになることにも注意してください。最後の方法は，\textbf{サーバの利用}です。これは環境を構築するのではなく，既に誰かがどこかで構築してくれたサービスを受け取るという，受動的な方法です。サービスを受け取ることのメリットは作り手の苦労を知る必要がないことですが，サービスが打ち切られると手元に何も残らなくなるという問題があります。この講義の受講者は，第一，第二の方法でどうしてもうまくいかなかった場合に，私の方からサーバ環境によるサービスを提供しますので，最後の砦だと思っておいてください\footnote{ケチくさいこと言わずに全員に提供しろよー，と思うかもしれませんが，サーバが貧弱で全受講生が一斉にアクセスするとサーバが止まってしまう可能性が高いからです。ちなみにサーバを強くするためにはたくさんのお金がかかりますので，私個人のサービスでは限界があるのです。}。

\section{導入方法1；ローカル環境での構築}
\section{導入方法2；仮想環境での構築}
\section{導入方法3；サーバの利用}
\end{document} 
